\section{Externe APIs und die Entscheidung für die TMDB API}
\setauthor{Markus Tran}

Ein zentraler Bestandteil der vorliegenden Anwendung ist die Darstellung von
Filminformationen und Streaming-Anbietern. Für die Bereitstellung dieser Daten
wurde eine externe API eingesetzt. Der folgende Abschnitt erläutert, was eine
API ist, warum eine externe Lösung gegenüber einer eigenen Datenbankpflege
bevorzugt wurde und nach welchen Kriterien die Entscheidung für die TMDB API
getroffen wurde.


\subsection{Definition und Funktion einer API}

Der Begriff API steht für \textit{Application Programming Interface} und
bezeichnet eine Schnittstelle, über die eine Anwendung auf Daten oder
Funktionalitäten eines externen Systems zugreifen kann. Eine API definiert
dabei, welche Anfragen gestellt werden können, in welchem Format die Antworten
zurückgegeben werden und wie die Authentifizierung gegenüber dem System erfolgt.

Im Kontext von Webanwendungen werden APIs häufig als REST-Schnittstellen
realisiert, die über HTTP-Anfragen angesprochen werden und ihre Antworten im
JSON-Format zurückgeben. Dieses Prinzip findet sich auch in der vorliegenden
Anwendung selbst wieder: Das Quarkus-Backend stellt dem Angular-Frontend
ebenfalls eine REST-API bereit und agiert damit gegenüber dem Frontend in
derselben Rolle wie eine externe API gegenüber dem Backend. Der Unterschied
liegt lediglich darin, dass das eigene Backend vollständig unter der Kontrolle
des Entwicklungsteams liegt, während eine externe API von einem Drittanbieter
betrieben und gepflegt wird.


\subsection{Eigene Datenhaltung vs. Auslagerung an eine externe API}

Bei der Konzeption der Anwendung stellte sich grundlegend die Frage, ob
Filmdaten eigenständig in der Datenbank verwaltet oder über eine externe
API bezogen werden sollen. Beide Ansätze haben spezifische Vor- und Nachteile,
die im Folgenden gegenübergestellt werden.

\subsubsection*{Eigene Datenhaltung}

Bei einer vollständig eigenständigen Datenhaltung würden alle Filmdaten --
Titel, Beschreibungen, Erscheinungsdaten, Bewertungen, Poster und
Streaming-Verfügbarkeiten -- in der eigenen Datenbank gespeichert und
gepflegt werden. Dies hätte den Vorteil, dass keinerlei Abhängigkeit von
einem externen Anbieter besteht, keine API-Schlüssel mit limitierten
Aufrufkontingenten benötigt werden und alle Daten vollständig kontrolliert
werden können.

Dem stehen jedoch erhebliche Nachteile gegenüber. Der globale Filmbestand
umfasst hunderttausende Einträge, was zu einer enormen Datenmenge führt.
Hinzu kommt der kontinuierliche Wartungsaufwand: Neue Filme müssten manuell
oder über eigene Crawler erfasst, Streaming-Verfügbarkeiten laufend
aktualisiert und veraltete Daten bereinigt werden. Dieser Aufwand übersteigt
den Rahmen eines Diplomarbeitsprojekts bei weitem und würde den Fokus von
der eigentlichen Applikationsentwicklung ablenken.

\subsubsection*{Auslagerung an eine externe API}

Die Auslagerung der Filmdaten an eine spezialisierte externe API vermeidet
die genannten Nachteile. Der Datenbankbetrieb, die Pflege des Filmbestands
sowie die Aktualisierung von Streaming-Verfügbarkeiten werden vollständig
vom API-Anbieter übernommen. Die eigene Anwendung konsumiert lediglich die
bereitgestellten Daten und kann sich auf ihre Kernfunktionalität
konzentrieren.

Als Nachteil ist die Abhängigkeit vom externen Anbieter zu nennen: Bei
Änderungen an der API-Struktur oder bei einer Einstellung des Dienstes
wäre eine Anpassung der Anwendung erforderlich. Zudem sind die Aufrufkontingente
je nach Tarifmodell begrenzt. Für den Anwendungsfall des vorliegenden Projekts
überwiegen jedoch die Vorteile der Auslagerung deutlich, weshalb diese Lösung
gewählt wurde.

\begin{table}[h]
\centering
\begin{tabular}{|p{3.5cm}|p{4.5cm}|p{4.5cm}|}
\hline
\textbf{Kriterium} & \textbf{Eigene Datenhaltung} & \textbf{Externe API} \\
\hline
Datenkontrolle & Vollständig & Abhängig vom Anbieter \\
Wartungsaufwand & Sehr hoch & Entfällt \\
Aktualität & Manuell sicherzustellen & Automatisch durch Anbieter \\
Datenmenge & Sehr groß & Nicht relevant \\
API-Kontingent & Nicht relevant & Begrenzt je nach Tarif \\
Implementierungs- & Sehr hoch & Gering \\
Aufwand & & \\
\hline
\end{tabular}
\caption{Vergleich: Eigene Datenhaltung vs. externe API}
\label{tab:api_vergleich}
\end{table}


\subsection{Auswahlkriterien für eine Film-API}

Für die Auswahl einer geeigneten Film-API wurden mehrere Kriterien definiert,
die für die Anforderungen der vorliegenden Anwendung relevant sind:

\begin{itemize}
    \item \textbf{Verfügbarkeit und Zugänglichkeit}: Die API soll über einen
    kostenlosen Zugang verfügen, der für Entwicklungs- und Testzwecke
    ausreichende Aufrufkontingente bietet.
    \item \textbf{Aktualität des Filmbestands}: Neue Filme sollen zeitnah
    in der Datenbank erfasst werden, sodass die Anwendung stets aktuelle
    Inhalte anzeigen kann.
    \item \textbf{Streaming-Provider-Integration}: Die API soll nicht nur
    Filmdaten, sondern auch Informationen zu Streaming-Verfügbarkeiten
    bereitstellen -- also auf welchen Plattformen ein Film verfügbar ist.
    Dies ist eine Kernanforderung der Anwendung.
    \item \textbf{Provider-Gesamtliste}: Zusätzlich zur filmspezifischen
    Verfügbarkeit soll die API eine vollständige Liste aller bekannten
    Streaming-Anbieter bereitstellen, die in der eigenen Datenbank
    synchronisiert werden kann.
    \item \textbf{Dokumentationsqualität}: Eine verständliche und vollständige
    API-Dokumentation ist Voraussetzung für eine effiziente Integration.
    \item \textbf{Verbreitung und Community}: Eine weit verbreitete API
    bietet erfahrungsgemäß eine stabilere Schnittstelle, regelmäßige
    Updates und eine aktive Entwickler-Community.
\end{itemize}

Im Rahmen der Evaluierung wurden unter anderem die \textit{Open Movie Database
API (OMDb)} sowie die \textit{TMDB API (The Movie Database)} betrachtet.
Die OMDb API bietet grundlegende Filmdaten, verfügt jedoch über keine
integrierte Bereitstellung von Streaming-Verfügbarkeiten, was eine
Kernanforderung der Anwendung darstellt. Damit schied diese Option aus.


\subsection{Entscheidung für die TMDB API}

Die Entscheidung fiel auf die \textbf{TMDB API (The Movie Database)}, da
sie alle definierten Auswahlkriterien erfüllt. Die TMDB ist eine der größten
öffentlich zugänglichen Film- und Seriendatenbanken und wird von einer
aktiven Community kontinuierlich gepflegt. Der kostenlose API-Zugang bietet
ausreichende Aufrufkontingente für den Betrieb der Anwendung.

Für die vorliegende Anwendung sind insbesondere folgende Endpunkte der
TMDB API relevant:

\begin{itemize}
    \item \texttt{/movie/\{id\}} -- Abrufen von Filmdaten anhand einer ID
    \item \texttt{/search/movie} -- Suche nach Filmen anhand eines Titels
    \item \texttt{/trending/movie/\{timewindow\}} -- Abrufen trendender Filme
    \item \texttt{/movie/\{id\}/watch/providers} -- Streaming-Anbieter
    eines Films, gefiltert nach Land
    \item \texttt{/watch/providers/movie} -- Vollständige Liste aller
    bekannten Streaming-Anbieter
\end{itemize}

Besonders hervorzuheben ist die Verfügbarkeit des
\texttt{/watch/providers}-Endpunkts, der nicht nur angibt, auf welchen
Plattformen ein Film verfügbar ist, sondern auch zwischen Flatrate-,
Leih- und Kaufangeboten unterscheidet. In der vorliegenden Anwendung
werden ausschließlich Flatrate-Anbieter berücksichtigt, da der Fokus
auf der Verwaltung aktiver Abonnements liegt. Die Authentifizierung
gegenüber der TMDB API erfolgt über einen Bearer-Token, der über den
Quarkus-Konfigurationsmechanismus mittels \texttt{@ConfigProperty}
sicher injiziert wird und nicht im Quellcode hinterlegt ist.